% header.tex — longtable-хук + xurl + microtype + инфраструктура глиф-фолбэков
\usepackage{graphicx}  % \includegraphics в raw-latex блоке фигур (build_preprint.py);
                       % pandoc сам его не тянет, т.к. markdown-картинок в доке нет
\usepackage{float}     % figure[H] -- картинка с подписью строго на месте, без float-страниц
\usepackage{xurl}
\usepackage{microtype}
\usepackage[htt]{hyphenat}
\usepackage{etoolbox}
\AtBeginEnvironment{longtable}{\scriptsize\setlength{\tabcolsep}{2.5pt}}
\setlength{\emergencystretch}{4em}
\usepackage{needspace}
% enotez -- в отличие от endnotes умеет двусторонние гиперссылки: номер в
% тексте кликабелен и ведёт на сноску внизу, а у сноски внизу есть обратная
% стрелка на место цитирования в тексте (backref) -- как в википедии.
% hyperfootnotes должен быть передан hyperref ДО его загрузки шаблоном
% pandoc, поэтому PassOptionsToPackage, а не прямой usepackage[...]{hyperref}.
\PassOptionsToPackage{hyperfootnotes=true}{hyperref}
\usepackage{multicol}
\usepackage{enotez}
\newcommand{\enotezNoHeading}[1]{}
\setenotez{backref=true, list-heading=\enotezNoHeading}
\let\footnote=\endnote
\renewcommand{\_}{\textunderscore\hspace{0pt}}
\usepackage{newunicodechar}
\newfontfamily\glyphfallback{TeX Gyre DejaVu Math}
% Если xelatex ругается на "Missing character" для конкретного юникод-символа
% (см. QC-проверку missing_chars в build_preprint.py) — добавить сюда:
% \newunicodechar{<символ>}{{\glyphfallback <символ>}}
